\chapter{Literature Review}
\label{chap:literature_review}

% Overview of the intersection congestion problem and the baseline study for Thapathali
\noindent The literature review was used to answer two practical questions for
this project. The first was which conventional signal-control method
should be used as a baseline for Thapathali. The second was which parts
of reinforcement learning are useful for a single-junction DQN
implementation. This framing matters because the project should not
simply attach a neural network to a traffic simulation. It needs a
baseline, a state representation that describes the local traffic
condition, and a reward function that matches the metrics used later in
evaluation.

The local Thapathali case study provides the field context for that
decision. It reports observed peak-hour volumes, delay-related
indicators, queue behaviour, and level of service for the existing
operation \cite{maharjan2023performance}.

% Fixed-time, actuated, and rule-based approaches to intersection signal timing
\section{Conventional Traffic Signal Control Methods}
\label{sec:conventional_methods}
\noindent For an isolated junction, traditional control usually starts with a
phase sequence, cycle length, and green split. Webster-type timing is a
common example. It uses traffic volume and simplifying assumptions to
produce a fixed-time plan for one intersection \cite{wei2020survey}.
These methods are still useful because they are easy to inspect and
implement. That is why fixed-time control is used as the baseline in this
report.

The weakness is the assumption behind the timing plan. A fixed plan
expects the design demand to remain representative during operation. At
Thapathali, where approach demand changes across the day, that assumption
is often too strong.

Actuated and rule-based controllers improve on fixed timing by using
detector inputs and predefined switching rules \cite{wei2020survey}.
Selection-based adaptive control chooses among prepared timing plans.
Optimization-based adaptive control optimizes an objective under a
simplified traffic-flow model \cite{wei2020survey}. These methods remain
valuable because they are interpretable and do not require long training
runs. Their limitation is calibration: if the real traffic state does not
match the simplified model or prepared plan, performance can drop.

Max-pressure control is especially relevant here because it uses queue
imbalance between upstream and downstream links to decide which phase
should receive service \cite{wei2020survey}. It is stronger than a simple
fixed-time baseline because it reacts to the observed queue state. Its
limitation is that the decision is still myopic: the controller reacts to
current pressure but does not learn a policy from repeated interaction.
For that reason, this project uses max-pressure as reward guidance rather
than as the complete controller.

% Framing traffic signal control as a sequential decision-making problem using reinforcement learning
\section{RL for Traffic Signal Control}
\label{sec:rl_traffic_control}
\noindent \gls{rl} changes the framing of traffic control from one-time timing calculation to sequential decision-making. The agent observes the traffic state, takes an action such as holding or changing the phase, and receives a reward that reflects the quality of the decision. This is useful because the correct action is not known in advance for every possible queue pattern. Instead of manually labelling the best signal decision for thousands of situations, the agent can learn through repeated simulation episodes \cite{wei2020survey}.

For a single-intersection project, the key design choices are the
observed state, allowed actions, and reward. Prior work commonly uses
lane-level counts, queue lengths, waiting time, and current signal phase
as state information. Rewards are usually based on delay, queue length,
pressure, or throughput \cite{wei2019colight,wei2020survey}. This shaped
the present design. PCU-weighted incoming loads, stopline waiting times,
outgoing traffic load, and active phase information were included because
each one describes a different control concern at Thapathali.

% Application of deep neural networks to approximate Q-values for traffic control decisions
\section{DRL and Value-Based Control with DQN}
\label{sec:drl_dqn_control}
\noindent \gls{drl} extends \gls{rl} by using neural networks to approximate value functions or policies, allowing the controller to learn from richer traffic states. The \gls{dqn} method is a practical value-based choice when the action space is discrete, because it estimates the value of each possible action and can be trained with experience replay to reduce instability \cite{wei2019presslight}. This fits the present project because the action space is intentionally limited to STAY and NEXT.

A recurring issue in \gls{drl} traffic-signal studies is that the
algorithm alone does not guarantee good control. Poor state choices can
hide important traffic conditions. Poorly shaped rewards can also push
the agent toward behaviour that improves one metric while damaging
another.

PressLight addresses this problem by linking \gls{drl} design with
max-pressure theory: the reward is shaped around pressure, and the state
is selected to support that objective \cite{wei2019presslight}. This is
the reason the present project uses a Max-Pressure based reward with an
added waiting-time penalty rather than a purely delay-based or
throughput-only reward.

Evaluation practice in \gls{drl} traffic control also influenced this report. Learned policies are normally compared against non-learning baselines such as fixed-time plans or transportation-inspired controllers, and the comparison is made using operational metrics such as travel time, delay, queue length, and throughput \cite{wei2019presslight,wei2020survey}. For this project, that means the DQN result is only meaningful if it is tested on the same Thapathali demand as the fixed signal plan and judged by measurable traffic indicators.

% Multi-intersection coordination and its relevance to future extensions of this work
\section{Network-Level Coordination as Context for Future Extension}
\label{sec:network_coordination}
\noindent The present report is limited to one intersection, but network-level studies are still useful for understanding future extensions. CoLight models traffic signal control as a cooperative multi-agent problem, where each intersection has its own agent and learns how much attention to give to neighbouring intersections under different traffic conditions \cite{wei2019colight}. Its observations include lane-level vehicle counts and current phase information, which are also used in this project, although without the neighbour-attention component.

The useful lesson from this work is not that the present minor project should immediately implement multi-agent control. Instead, it shows why the state design should not be tied too tightly to one hand-written scenario. If the Thapathali controller is later extended to nearby junctions, lane-level features and explicit phase encoding will be easier to reuse than a highly specific set of manually engineered rules.

% Phase-level action representation and exploiting geometric symmetry for sample efficiency
\section{Structured Policy Learning via Phase Competition and Symmetry}
\label{sec:structured_policy_learning}
\noindent Another relevant direction is structured policy learning, where the controller is designed around meaningful traffic phases rather than arbitrary low-level signal changes. A phase-competition approach argues that many traffic states are equivalent under rotation or flipping, and that an \gls{rl} agent should generalise across such symmetric cases instead of relearning each one separately \cite{zheng2019learning}. The method models competition among phases and gives priority to movements with higher demand when conflicts exist \cite{zheng2019learning}.

For this project, the practical takeaway is to keep the action space small and tied to real phase movement. Allowing the agent to control every signal lamp independently would create a larger exploration problem and could produce unsafe or unrealistic phase combinations. The STAY/NEXT action design is simpler, but it respects the existing phase cycle and keeps the learning problem manageable.

% Role of SUMO simulation and local data in validating the proposed approach
\section{Simulation and Realistic Data Grounding}
\label{sec:simulation_grounding}
\noindent Simulation is the practical bridge between the reviewed methods and this project. Real-world testing at Thapathali would require hardware access, traffic authority approval, and safety safeguards that are outside the scope of a minor project. \gls{sumo} is suitable here because it can simulate vehicle movement, expose lane-level measurements, and accept signal-control actions during runtime \cite{zheng2019learning}.

The local Thapathali data is equally important. The case study reports peak-hour traffic volumes in \gls{pcu}, approach-level distributions, degree of saturation, average delay, level of service, and queue-related measures \cite{maharjan2023performance}. These values make the simulation less arbitrary: they provide the demand pattern used to generate vehicle flows and the operational context against which the result can be interpreted. The limitation is that the available OD proportions are not available for every interval, so some interpolation and allocation assumptions are still required.

% Synthesis of reviewed literature and positioning of this project relative to the state of the art
\section{Summary and Project Positioning}
\label{sec:lit_review_summary}
\noindent The reviewed work suggests a clear direction for this project. Fixed-time and rule-based controllers are still necessary baselines because they are understandable and widely used. Max-pressure control provides a stronger traffic-theory reference, especially for reward design. \gls{drl} methods such as \gls{dqn} make adaptive control possible, but the literature also shows that careless state, action, or reward choices can make learning unstable or difficult to interpret.

Based on this review, the project adopts a deliberately limited design: a single-intersection \gls{dqn} controller trained in \gls{sumo}, using local Thapathali demand data, a 21-dimensional traffic state, two discrete actions, and a Max-Pressure based reward. This does not cover every advanced direction in the literature, such as attention-based multi-agent coordination or symmetry-aware phase learning, but it uses the parts that are feasible and relevant for the present project scope \cite{wei2020survey,wei2019presslight,wei2019colight,zheng2019learning,maharjan2023performance}.

\begin{table}[H]
    \centering
    \small
    \caption{Comparison of Traffic Signal Control Methods}
    \label{tab:traffic_control_methods}
    \begin{tabular}{p{2.5cm}p{4cm}p{4cm}p{4.5cm}}
        \toprule
        \textbf{Method} & \textbf{Control Logic} & \textbf{Adaptivity} & \textbf{Data Needs} \\
        \midrule
        Fixed Time & Pre-set cycle length, phase order, and green splits (static plan). & Low: no real-time response; only changes if manually retimed. & Historical/observed volumes for timing design; no training required. \\
        \addlinespace
        Max Pressure & Selects phase with highest pressure using upstream vs downstream queue imbalance (rule-based). & High: reacts instantly to queues, but remains a fixed greedy rule (no learning). & Real-time queue/occupancy (ideally downstream too); no training dataset. \\
        \addlinespace
        DQN-based DRL & Learns a signal policy by maximizing reward through simulation interaction. & High: learns and adapts a policy to objectives and constraints. & Training episodes in \gls{sumo}; state features; compute and tuning effort. \\
        \bottomrule
    \end{tabular}
\end{table}
